\documentclass[11pt]{article}
\usepackage[a4paper,margin=23mm]{geometry}
\usepackage{amsmath,amssymb,amsthm,mathtools,bm}
\usepackage{slashed}
\usepackage{booktabs,longtable,array}
\usepackage{enumitem}
\usepackage{xcolor}
\usepackage{microtype}
\usepackage[hidelinks]{hyperref}

\setlength{\emergencystretch}{3em}
\newtheorem{theorem}{Theorem}[section]
\newtheorem{proposition}[theorem]{Proposition}
\newtheorem{lemma}[theorem]{Lemma}
\newtheorem{corollary}[theorem]{Corollary}
\theoremstyle{definition}
\newtheorem{definition}[theorem]{Definition}
\newtheorem{gate}[theorem]{Fail-closed gate}
\theoremstyle{remark}
\newtheorem{remark}[theorem]{Remark}

\newcommand{\dd}{\mathrm d}
\newcommand{\ii}{\mathrm i}
\newcommand{\ZZ}{\mathbb Z}
\newcommand{\RR}{\mathbb R}
\newcommand{\CC}{\mathbb C}
\newcommand{\Tr}{\operatorname{Tr}}
\newcommand{\Arf}{\operatorname{Arf}}
\newcommand{\ind}{\operatorname{ind}}
\newcommand{\SF}{\operatorname{SF}}
\newcommand{\Pf}{\operatorname{Pf}}
\newcommand{\status}[1]{\textcolor{blue!60!black}{\textsf{#1}}}
\newcolumntype{L}[1]{>{\raggedright\arraybackslash}p{#1}}

\title{AP-E3 APS Generators and Semisimple Global Forms:\\
Product/Defect Mod-Two Indices, Regulator Non-Uniqueness,\\
and a Simply Connected \(Sp(4)\) Candidate}
\author{Route F independent research note}
\date{16 July 2026}

\begin{document}
\maketitle

\begin{abstract}
We evaluate the two reduced generators of
\(\Omega_4^{\mathrm{Spin}}(S^3\times S^2)\) in explicit solvable fermionic
models and systematically search semisimple global forms that evade the
previous \(SU(3)\) charge obstruction.  The generators are
\[
 G_3=(S^1_{\rm R}\times S^3,\operatorname{pr}_{S^3},{\rm pt}),
 \qquad
 G_2=(T^2_{\rm RR}\times S^2,{\rm pt},\operatorname{pr}_{S^2}).
\]
The periodic-circle real Dirac operator has one zero mode and the odd torus
chiral Dirac operator has mod-two index one.  A codimension-three defect
regulator therefore evaluates \((-1)^{\nu_3}\), while an Arf-decorated
codimension-two regulator evaluates \((-1)^{\nu_2}\).  Stacking these two
invertible regulators produces all four character values on \((G_3,G_2)\).
By contrast, the untwisted ambient product Dirac spectra are chirality paired
and provide only a reference phase \(+1\) in a factorized regulator.

This is a controlled negative result: it computes a product/defect mod-two
index character table, not the unique eta phase of the charged two-colour UV
theory.  The current data do not specify the complete Euclidean
Dirac--Yukawa operator, Pauli--Villars mass signs, or five-dimensional
mapping-torus APS problem.  Heavy invertible spectators can independently
flip both torsion signs without changing the degree-five differential
curvature.  Hence neither sign is fixed from the present effective action.

For global forms, a short theorem shows that any unbroken
\([SU(2)_c\times U(1)]/\ZZ_2\) excludes a charge-one colour singlet, although
it admits \(\bm2_{+1}\).  Thus the diagonal quotient itself, not merely
\(SU(3)\) branching, is the obstruction.  Three unquotiented semisimple
forms pass the charge screen.  The strongest simple candidate in the scanned
set is the simply connected group
\(Sp(4)\equiv USp(4)\simeq Spin(5)\), with
\[
 \bm5\to\bm2_{+1}\oplus\bm2_{-1}\oplus\bm1_0,
 \qquad
 \bm4\to\bm2_0\oplus\bm1_{+1}\oplus\bm1_{-1}.
\]
Two Dirac \(\bm5\)'s and two complex scalar \(\bm4\)'s give the required
two-flavour quarks and an exact copy-space \(SU(2)_\phi\) doublet.  One real
adjoint can break to \(SU(2)_c\times U(1)\) and split all displayed partners
at tree level.  The one-loop coefficient is \(b_0^{Sp(4)}=15/2>0\), and the
displayed vectorlike spectrum passes local and Witten-parity checks.
Nevertheless, the massless charged component is tuned, the heavy eta
threshold is unmatched, and no \(Sp(4)\) strong-phase or soliton calculation
is supplied.  The exactly-two operator is therefore realized only at the
group-theory level; the exactly-two gauge/unit-cell rule is not rederived.
No Route-E or degree-one-portal promotion is made.
\end{abstract}

\tableofcontents

\section{Scope and claim boundary}

The previous global audit separated the unique ordinary degree-five
differential character from two spin-torsion characters.  With
\(X=S^3\times S^2\), stable splitting gives
\begin{equation}
 \Omega_4^{\mathrm{Spin}}(X)
 \simeq \ZZ\oplus\ZZ_2\oplus\ZZ_2,
 \qquad
 \widetilde\Omega_4^{\mathrm{Spin}}(X)
 \simeq\ZZ_2\oplus\ZZ_2.
 \label{eq:bordism}
\end{equation}
After the field-independent gravitational factor is removed, the general
spin refinement has the form
\begin{equation}
 Z_{n,\epsilon_3,\epsilon_2}
 =Z_n^{\rm diff}
 (-1)^{\epsilon_3\nu_3+\epsilon_2\nu_2},
 \qquad (\epsilon_3,\epsilon_2)\in\ZZ_2^2.
 \label{eq:four-characters}
\end{equation}
The purpose of this note is to answer two narrower questions:
\begin{enumerate}[label=(Q\arabic*)]
\item what can be computed on explicit representatives of the two reduced
  generators with a fully stated product/defect regulator;
\item which low-rank semisimple global forms contain both
  \(q\sim\bm2_{+1}\) and \(\phi\sim\bm1_{+1}\), so that
  \(qq(\phi^\dagger)^2\) is neutral.
\end{enumerate}

Three distinctions will be maintained throughout.
\begin{enumerate}
\item A paired spectrum of an untwisted four-dimensional product Dirac
  operator is a \emph{reference spectral control}; it is not, by itself, a
  five-dimensional Dai--Freed computation.
\item A defect mod-two index computes the phase of the explicitly selected
  invertible regulator.  It proves that a torsion character is realizable,
  not that charged QC2D selects it.
\item Charge neutrality and branching prove an operator-level selection
  rule.  They do not derive the microscopic exactly-two gauge/unit-cell
  constraint, threshold naturalness, or the nonperturbative phase.
\end{enumerate}
This discipline follows the APS/Dai--Freed treatment of determinant phases
\cite{AtiyahPatodiSinger1975,DaiFreed1994,Yonekura2019} and the bordism
approach to four-dimensional anomalies
\cite{DavighiGripaiosLohitsiri2020,Saito2024}.

\section{Explicit representatives of both reduced generators}

\subsection{Pontryagin--Thom invariants}

For a sigma-model field \(\Phi=(g,s):M_4\to S^3\times S^2\), choose regular
values \(p_3\in S^3\), \(p_2\in S^2\).  Their transverse inverse images are
\begin{equation}
 L_g=g^{-1}(p_3),\qquad \dim L_g=1,
 \qquad
 \Sigma_s=s^{-1}(p_2),\qquad \dim\Sigma_s=2.
 \label{eq:inverse-images}
\end{equation}
The target orientations frame the normal bundles.  Together with the spin
structure on \(M_4\), they induce spin structures on \(L_g\) and
\(\Sigma_s\).  The reduced invariants are
\begin{equation}
 \nu_3=[L_g]\in\Omega_1^{\mathrm{Spin}}\simeq\ZZ_2,
 \qquad
 \nu_2=\Arf(\Sigma_s)\in\Omega_2^{\mathrm{Spin}}\simeq\ZZ_2.
 \label{eq:nu-definitions}
\end{equation}

Take unit-radius factors and define
\begin{align}
 G_3&:=
 (S^1_{\rm R}\times S^3,
  g=\operatorname{pr}_{S^3},s={\rm const}),
 \label{eq:G3}\\
 G_2&:=
 (T^2_{\rm RR}\times S^2,
  g={\rm const},s=\operatorname{pr}_{S^2}).
 \label{eq:G2}
\end{align}
Here R denotes periodic spin structure.  The inverse images are
\begin{equation}
 L_g=S^1_{\rm R}\times\{p_3\},
 \qquad
 \Sigma_s=T^2_{\rm RR}\times\{p_2\}.
\end{equation}
The periodic circle is the nonbounding generator of
\(\Omega_1^{\mathrm{Spin}}\), while the RR torus is the odd spin surface with
Arf invariant one \cite{Atiyah1971}.  Therefore
\begin{equation}
 (\nu_3(G_3),\nu_2(G_3))=(1,0),
 \qquad
 (\nu_3(G_2),\nu_2(G_2))=(0,1).
 \label{eq:generator-values}
\end{equation}

\subsection{Circle spectrum and mod-two spectral flow}

Let the circle have circumference \(2\pi\).  Its real Dirac spectrum is
\begin{equation}
 D_{S^1,\alpha}=-\ii\frac{\dd}{\dd x},
 \qquad
 \lambda_n=n+\alpha,
 \qquad n\in\ZZ,
 \label{eq:circle-spectrum}
\end{equation}
where \(\alpha=0\) for R and \(\alpha=1/2\) for NS spin structure.  Hence
\begin{equation}
 \dim_{\RR}\ker D_{S^1,0}=1\pmod2,
 \qquad
 \ker D_{S^1,1/2}=0.
 \label{eq:circle-ko-index}
\end{equation}
For the explicit mass path
\begin{equation}
 H_\alpha(t)=D_{S^1,\alpha}+t,
 \qquad -\frac14\le t\le\frac14,
 \label{eq:circle-mass-path}
\end{equation}
only the \(n=0\) R eigenvalue crosses zero.  Thus
\begin{equation}
 \SF_2(H_0)=1,
 \qquad
 \SF_2(H_{1/2})=0,
 \label{eq:circle-sf}
\end{equation}
and APS identifies this with the mod-two index of
\(\partial_t+H_\alpha(t)\).  A real Pfaffian changes sign exactly in the R
case.

\subsection{Odd-torus chiral index and Arf phase}

On a square torus of side \(2\pi\), the chiral Dirac eigenvalues may be
normalized as
\begin{equation}
 D^+_{\alpha\beta}:\quad
 \lambda_{mn}=(m+\alpha)+\ii(n+\beta),
 \qquad (m,n)\in\ZZ^2.
 \label{eq:torus-spectrum}
\end{equation}
Only \((\alpha,\beta)=(0,0)\) has a zero mode.  Consequently
\begin{equation}
 \ind_2 D^+_{T^2_{\rm RR}}
 =\dim_{\CC}\ker D^+_{00}\pmod2=1,
 \qquad
 \ind_2D^+_{\alpha\beta}=0
 \quad\text{otherwise}.
 \label{eq:torus-index}
\end{equation}
The equality
\begin{equation}
 \Arf(T^2,\alpha,\beta)=\ind_2D^+_{\alpha\beta}
 \label{eq:arf-index}
\end{equation}
identifies the Majorana Pfaffian sign on \(G_2\).  This is a KO/mod-two
index statement and needs no continuum extrapolation.

\section{Reference ambient spectra versus the actual APS problem}

\subsection{Paired product spectra}

The ordinary Dirac spectra on unit spheres are
\begin{align}
 \operatorname{spec}(D_{S^3})
 &=\left\{\pm(k+\tfrac32),\quad
   \operatorname{mult}=(k+1)(k+2)\right\}_{k\ge0},
 \label{eq:s3-spectrum}\\
 \operatorname{spec}(D_{S^2})
 &=\left\{\pm(\ell+1),\quad
   \operatorname{mult}=2(\ell+1)\right\}_{\ell\ge0}.
 \label{eq:s2-spectrum}
\end{align}
For product metrics,
\begin{align}
 |\lambda(D_{S^1_{\rm R}\times S^3})|^2
 &=n^2+(k+\tfrac32)^2\ge\frac94,
 \label{eq:s1s3-gap}\\
 |\lambda(D_{T^2_{\rm RR}\times S^2})|^2
 &=m^2+n^2+(\ell+1)^2\ge1.
 \label{eq:t2s2-gap}
\end{align}
Both ambient four-dimensional operators are gapped, and chirality pairs
every nonzero eigenvalue with its negative.  A symmetric, factorized
Pauli--Villars convention therefore assigns reference phase \(+1\) to both
products.

\begin{remark}[What the paired spectrum does not prove]
The eta invariant controlling a four-dimensional microscopic anomaly belongs
to a specified odd-dimensional mapping-torus or APS problem.  Equations
\eqref{eq:s1s3-gap}--\eqref{eq:t2s2-gap} do not determine that problem.  They
only verify that an untwisted ambient product determinant has no unpaired
four-dimensional eigenvalue.  Calling this calculation ``the charged-QC2D
eta invariant'' would be incorrect.
\end{remark}

\subsection{An explicit defect Dirac--Yukawa family}

We now construct a separate, computable regulator family.  In a tubular
neighbourhood of a transverse inverse image \(N=\Phi^{-1}(p)\), use oriented
normal coordinates \(u^a\), \(a=1,\dots,r\), and an auxiliary real Clifford
module with matrices \(\Gamma_a\).  A local Dirac--Yukawa operator is
\begin{equation}
 \mathcal D_{r,\mu}
 =\slashed D_{M_4}\otimes\bm1
 +\mu\sum_{a=1}^{r}u^a(\Phi)\Gamma_a.
 \label{eq:defect-yukawa}
\end{equation}
Using a cutoff in the tubular neighbourhood and a positive Pauli--Villars
mass outside gives a global Fredholm family.  The normal harmonic-oscillator
modes pair except for a single localized mode on \(N\), as in the
Jackiw--Rebbi mechanism \cite{JackiwRebbi1976}.  Gapping that localized mode
by the nontrivial one- or two-dimensional invertible spin phase gives
\begin{align}
 R_3[M_4,g]&=(-1)^{\ind_2D_{L_g}}
            =(-1)^{\nu_3(M_4,g)},
 \label{eq:R3}\\
 R_2[M_4,s]&=(-1)^{\ind_2D^+_{\Sigma_s}}
            =(-1)^{\Arf(\Sigma_s)}
            =(-1)^{\nu_2(M_4,s)}.
 \label{eq:R2}
\end{align}
The full APS operator factorizes in the large-\(\mu\) limit into paired
normal modes and the defect operator.  Therefore its mod-two index is the
right-hand side of Eqs.~\eqref{eq:R3}--\eqref{eq:R2}.  This establishes an
explicit regulator realization of each character.

The resulting character table is
\begin{center}
\begin{tabular}{cccc}
\toprule
\((\epsilon_3,\epsilon_2)\) & regulator stack & phase on \(G_3\) & phase on \(G_2\)\\
\midrule
\((0,0)\) & reference & \(+1\) & \(+1\)\\
\((1,0)\) & \(R_3\) & \(-1\) & \(+1\)\\
\((0,1)\) & \(R_2\) & \(+1\) & \(-1\)\\
\((1,1)\) & \(R_3R_2\) & \(-1\) & \(-1\)\\
\bottomrule
\end{tabular}
\end{center}

\begin{theorem}[Regulator non-uniqueness]
The differential curvature and all perturbative local terms held fixed do
not determine \((\epsilon_3,\epsilon_2)\).  Stacking either heavy invertible
defect regulator changes exactly one torsion sign.  Hence a microscopic UV
Dirac--Yukawa operator and its regulator are necessary data for selecting the
two signs.
\end{theorem}
\begin{proof}
Equations~\eqref{eq:R3} and \eqref{eq:R2} are homomorphisms from reduced spin
bordism to \(\{\pm1\}\).  Equation~\eqref{eq:generator-values} shows that
their evaluation matrix is the identity over \(\ZZ_2\), so they form a basis
of
\(\operatorname{Hom}(\ZZ_2\oplus\ZZ_2,U(1))\).  They are torsion phases and
have zero de Rham curvature.  Stacking therefore leaves the fixed
degree-five curvature unchanged while toggling either sign.
\end{proof}

\begin{gate}[Charged-QC2D determinant]
The product/defect construction proves the four possible phase assignments
and supplies a regulator-dependence negative control.  The current Route-F
files do not specify the complete charged-QC2D Euclidean operator
\begin{equation}
 \mathcal D_{\rm UV}
 =\gamma^\mu(\nabla_\mu+A_\mu)
 +P_L\mathcal M(g,s)+P_R\mathcal M(g,s)^\dagger,
 \label{eq:actual-uv-operator}
\end{equation}
including every heavy representation, Yukawa matrix, Pauli--Villars mass
sign, and mapping-torus extension.  Thus
\begin{equation}
 \begin{gathered}
 \texttt{epsilon\_3\_selected\_by\_charged\_qc2d\_uv=false},\\
 \texttt{epsilon\_2\_selected\_by\_charged\_qc2d\_uv=false}.
 \end{gathered}
 \label{eq:eta-fail-closed}
\end{equation}
\end{gate}

\section{A global-form theorem before branching}

Let
\begin{equation}
 H_0=SU(2)_c\times U(1)_H,
 \qquad
 H_1=\frac{SU(2)_c\times U(1)_H}
 {\langle(-\bm1,e^{\ii\pi})\rangle}.
 \label{eq:H-forms}
\end{equation}
Normalize \(U(1)_H\) so that its smallest desired charge is one.  A term
\((j,q)\) descends to \(H_1\) precisely when the quotient generator acts
trivially:
\begin{equation}
 (-1)^{2j}(-1)^q=1
 \quad\Longleftrightarrow\quad
 2j+q=0\pmod2.
 \label{eq:quotient-condition}
\end{equation}

\begin{theorem}[Diagonal-quotient obstruction]
The diagonal quotient \(H_1\) admits a colour doublet \(\bm2_{+1}\) but
excludes a colour singlet \(\bm1_{+1}\).  Therefore any semisimple embedding
with unbroken global form \(H_1\) fails the charge-one/exactly-two field
requirement independently of its Lie algebra branching.
\end{theorem}
\begin{proof}
For \(\bm2_{+1}\), \(2j+q=1+1=2\), so Eq.~\eqref{eq:quotient-condition}
holds.  For \(\bm1_{+1}\), it gives \(0+1=1\), so it fails.  Since this is a
statement about the kernel of the group homomorphism, no higher
representation can evade it.
\end{proof}

The exact operator charge condition is
\begin{equation}
 Q\!\left[qq(\phi^\dagger)^m\right]
 =2X_q-mX_\phi.
 \label{eq:dressing-charge}
\end{equation}
With \(X_q=X_\phi=1\), the first neutral power is \(m=2\).  Thus an
\emph{unquotiented} \(H_0\) solves the operator charge equation.  It does not
by itself prove that a microscopic unit cell has exactly two constituents.

\section{Systematic low-rank semisimple search}

\subsection{Baseline failures}

For
\(SU(3)\to[SU(2)_c\times U(1)]/\ZZ_2\),
\begin{equation}
 \bm3\to\bm2_{+1}\oplus\bm1_{-2}.
 \label{eq:su3-branch}
\end{equation}
The quotient theorem proves that no \(\bm1_{+1}\) occurs in any
\(SU(3)\) representation.  This recovers the previous no-go without relying
on a finite representation scan.

The same warning applies to covers and quotients.  The direct product
\(SU(2)_c\times SU(2)_H\) contains the required representations, whereas
\(SO(4)=[SU(2)_c\times SU(2)_H]/\ZZ_2\) removes the scalar.  Likewise the
spinor \(\bm4\) of \(Spin(5)\) does not descend to \(SO(5)\).  Global form is
therefore physical data, not notation.

\subsection{Minimal product candidate}

Take the unquotiented semisimple group
\begin{equation}
 G_{\rm prod}=SU(2)_c\times SU(2)_H
 \xrightarrow{\langle\Sigma_H\rangle}
 SU(2)_c\times U(1)_H,
 \label{eq:product-breaking}
\end{equation}
where a real \((\bm1,\bm3)\) scalar has
\(\langle\Sigma_H\rangle=V T_H^3\), and the doublet weights are
\(q=\pm1\).  Then
\begin{align}
 (\bm2,\bm2)&\to\bm2_{+1}\oplus\bm2_{-1},
 \label{eq:product-fermion}\\
 (\bm1,\bm2)&\to\bm1_{+1}\oplus\bm1_{-1}.
 \label{eq:product-scalar}
\end{align}
Two Dirac bifundamentals supply two low-energy quark flavours.  Two complex
scalar \((\bm1,\bm2)\) copies \(\Phi_I\), \(I=1,2\), retain an exact
copy-space \(SU(2)_\phi\) when their couplings are proportional to
\(\delta_{IJ}\).  The two desired charge-one components then form the
\(SU(2)_\phi\) doublet used in
\(\operatorname{Sym}^2\bm2=\bm3\).

The group is already a viable minimal semisimple backup.  Its limitation is
that each required field comes with an opposite-charge partner, and the
massless branch requires the threshold split discussed below.

\subsection{Best simple candidate in the scanned set: simply connected
\texorpdfstring{\(Sp(4)\)}{Sp(4)}}

We use physics notation
\begin{equation}
 Sp(4)\equiv USp(4)\simeq Spin(5),
 \label{eq:sp4-notation}
\end{equation}
whose fundamental has dimension four.  It contains
\(Spin(4)=SU(2)_c\times SU(2)_H\).  Under this subgroup
\begin{align}
 \bm4&\to(\bm2,\bm1)\oplus(\bm1,\bm2),
 \label{eq:sp4-4-spin4}\\
 \bm5&\to(\bm2,\bm2)\oplus(\bm1,\bm1),
 \label{eq:sp4-5-spin4}\\
 \bm{10}&\to(\bm3,\bm1)\oplus(\bm1,\bm3)
                    \oplus(\bm2,\bm2).
 \label{eq:sp4-10-spin4}
\end{align}
These formulas also follow directly from
\(\bm5=\wedge^2\bm4-\bm1\) and
\(\bm{10}=\operatorname{Sym}^2\bm4\)
\cite{Slansky1981}.  Breaking the second factor by a real adjoint
\(\Sigma\in\bm{10}\), with its vacuum in the
\((\bm1,\bm3)_0\) direction, gives
\begin{align}
 \bm4&\to\bm2_0\oplus\bm1_{+1}\oplus\bm1_{-1},
 \label{eq:sp4-4}\\
 \bm5&\to\bm2_{+1}\oplus\bm2_{-1}\oplus\bm1_0,
 \label{eq:sp4-5}\\
 \bm{10}&\to
 \bm1_{+2}\oplus\bm2_{+1}\oplus\bm3_0\oplus\bm1_0
 \oplus\bm2_{-1}\oplus\bm1_{-2}.
 \label{eq:sp4-10}
\end{align}

The subgroup embedding is injective.  In particular,
\((-\bm1_c,-\bm1_H)\) maps to the nontrivial central element of \(Spin(5)\),
not to the identity.  After \(SU(2)_H\to U(1)_H\), the unbroken global form
is therefore \(H_0\), so both odd-charge representations are legitimate.
If one instead quotients \(Spin(5)\) to \(SO(5)\), this central element is
identified with the identity and the spinor \(\bm4\) is lost.  The scalar
requirement forces the simply connected cover.

An explicit field list is
\begin{equation}
 \Psi_i\in\bm5\ \text{Dirac},\quad i=1,2;
 \qquad
 \Phi_I\in\bm4\ \text{complex},\quad I=1,2;
 \qquad
 \Sigma\in\bm{10}\ \text{real}.
 \label{eq:sp4-field-list}
\end{equation}
The two \(\Phi_I\) copies are essential.  Their copy index carries the exact
global \(SU(2)_\phi\); the two gauge weights inside a single \(\bm4\) are not
a substitute for this flavour doublet.  With copy-symmetric couplings, the
light \(\bm1_{+1}\) components are \(\phi_I\), and
\begin{equation}
 qq(\phi_I^\dagger\phi_J^\dagger)_{\rm sym}
 \in\operatorname{Sym}^2\bm2_\phi=\bm3_\phi.
 \label{eq:sp4-triplet}
\end{equation}
Equation~\eqref{eq:sp4-triplet} establishes the desired group representation
of the local operator.  It does not prove the exactly-two microscopic Gauss
rule.

\subsection{A secondary simple candidate: \texorpdfstring{\(SU(4)\)}{SU(4)}}

The block embedding
\begin{equation}
 \iota_4(U,z)=\operatorname{diag}(zU,z^{-1},z^{-1})
 \in SU(4)
 \label{eq:su4-embedding}
\end{equation}
has trivial kernel.  Consequently
\begin{align}
 \bm4&\to\bm2_{+1}\oplus2\,\bm1_{-1},
 \label{eq:su4-4}\\
 \overline{\bm4}&\to\bm2_{-1}\oplus2\,\bm1_{+1}.
 \label{eq:su4-bar4}
\end{align}
Two suitably oriented adjoint vacua can reduce the centralizer to the direct
\(SU(2)_c\times U(1)\) while preserving the block charges.  The fundamental
also has a useful uniqueness property.  For an \(SU(N)\) block embedding
with doublet charge \(a\) and each of the other \(N-2\) entries of charge
\(b\), tracelessness gives
\begin{equation}
 2a+(N-2)b=0.
 \label{eq:sun-traceless}
\end{equation}
Demanding equal-magnitude quark and conjugate-singlet charges, \(a=-b\ne0\),
forces
\begin{equation}
 (4-N)a=0\quad\Longrightarrow\quad N=4.
 \label{eq:su4-unique}
\end{equation}
Thus \(SU(4)\) is the unique member of this fundamental block family with
the desired equality.  Quotienting by \(\{-\bm1_4\}\) forbids both
\(\bm4\) and \(\overline{\bm4}\), so again the simply connected form is
required.  We rank this below \(Sp(4)\) because the required breaking and
copy-flavour protection use more independent ingredients.

\section{Tree-level partner decoupling}

\subsection{Universal adjoint split}

In all three passing candidates, the desired and conjugate components have
weights \(q=\pm1\) under the broken generator.  A vectorlike Dirac multiplet
admits
\begin{equation}
 \mathcal L_{\Psi}
 =-\overline\Psi(m+y\Sigma)\Psi,
 \qquad
 M_q=m+q\,yV.
 \label{eq:fermion-masses}
\end{equation}
Choosing
\begin{equation}
 m=-yV
 \label{eq:fermion-tuning}
\end{equation}
gives
\begin{equation}
 M_{+}=0,\qquad M_{-}=-2yV.
 \label{eq:fermion-split}
\end{equation}
For \(Sp(4)\), the neutral singlet in \(\bm5\) has mass \(M_0=-yV\).  Thus
all unwanted fermions are heavy at tree level when \(yV\ne0\).

For each scalar copy, take
\begin{equation}
 V_2=\Phi_I^\dagger(m_\Phi^2+\kappa\Sigma)\Phi_I.
 \label{eq:scalar-masses}
\end{equation}
The benchmark
\begin{equation}
 m_\Phi^2=1,\qquad \kappa V=-2
 \label{eq:scalar-benchmark}
\end{equation}
gives
\begin{equation}
 m^2_{\phi,+}=-1,\qquad
 m^2_{\phi,-}=3,
 \qquad
 m^2_{\phi,0}=1\quad(Sp(4)).
 \label{eq:scalar-split}
\end{equation}
Quartics can stabilize the desired charged vacuum while leaving the displayed
partners positive at tree level.

\begin{gate}[Radiative protection]
Equation~\eqref{eq:fermion-tuning} is a cancellation between independent
operators.  No symmetry protecting it has been specified.  Heavy gauge,
adjoint, and partner loops generically induce
\begin{equation}
 \delta M_+\sim\frac{g_G^2}{16\pi^2}\,yV
 \label{eq:threshold-estimate}
\end{equation}
up to representation- and scheme-dependent coefficients.  The candidate is
therefore an explicit tree-level spectrum, not a natural all-scale
decoupling theorem.
\end{gate}

\subsection{Orientation and the unmatched heavy determinant}

The retained colour doublet has \(X_q=+1\) and two colours.  Its light-field
mixed coefficient is therefore
\begin{equation}
 k_{\rm light}=N_cX_q=2(+1)=+2.
 \label{eq:light-k}
\end{equation}
This fixes the algebraic orientation convention.  Integrating out the
\(q=-1\) partner and neutral fields also generates regulator-dependent local
and torsion counterterms.  Until their determinant phases are matched across
the breaking threshold, Eq.~\eqref{eq:light-k} is not an all-scale
Dai--Freed derivation of \(k=+2\).

\section{Gauge anomalies and one-loop running}

\subsection{Displayed anomaly checks}

All fermions in the candidate lists are Dirac fields in complete deep-group
representations.  Their left- and right-handed perturbative gauge anomalies
cancel.  For the product group, two Dirac \((\bm2,\bm2)\)'s give
\begin{equation}
 N^{(c)}_{\rm Weyl\ doublet}
 =N^{(H)}_{\rm Weyl\ doublet}
 =N_f\,(2\ {\rm chiralities})\,(2\ {\rm spectator\ dimension})=8,
 \label{eq:product-witten}
\end{equation}
which is even.  In the \(Sp(4)\) candidate, the displayed Dirac fields double
every chiral mod-two phase.  After breaking, the two light Dirac colour
doublets give four left-Weyl colour doublets, also even.  These are the
ordinary Witten-parity checks \cite{Witten1982,JiaYi2024}.

\begin{gate}[Full gauge-bordism phase]
Vectorlike cancellation and the \(\pi_4\) parity test are necessary checks,
not a substitute for evaluating the complete
\(\Omega_5^{\mathrm{Spin}}(BSp(4))\) Dai--Freed character with every Higgs and
background bundle retained.  We therefore set
\begin{equation}
 \texttt{best\_candidate\_full\_dai\_freed\_gauge\_bordism\_audit\_complete=false}.
 \label{eq:gauge-bordism-open}
\end{equation}
This is deliberately conservative even though the displayed matter is
vectorlike.
\end{gate}

\subsection{Asymptotic-freedom screen}

For Dirac fermions, real scalars, and complex scalars, respectively, use
\begin{equation}
 b_0=\frac{11}{3}C_2(G)
 -\frac43\sum_{\rm Dirac}T(R_f)
 -\frac16\sum_{\rm real}T(R_s)
 -\frac13\sum_{\rm complex}T(R_s).
 \label{eq:beta-general}
\end{equation}
Our conventions are
\begin{equation}
 T_{SU(2)}(\bm2)=\frac12,\quad T_{SU(2)}(\bm3)=2,
 \qquad
 T_{Sp(4)}(\bm4)=\frac12,\quad T_{Sp(4)}(\bm5)=1,
 \quad T_{Sp(4)}(\bm{10})=3.
 \label{eq:indices}
\end{equation}
For the product candidate, the bifundamental spectator dimension makes
\(\sum_fT=2\) for each factor.  With one real \(SU(2)_H\) adjoint and two
complex \(H\) doublets,
\begin{align}
 b_{0,c}^{\rm prod}
 &=\frac{22}{3}-\frac43(2)=\frac{14}{3}>0,
 \label{eq:b0-product-c}\\
 b_{0,H}^{\rm prod}
 &=\frac{22}{3}-\frac43(2)-\frac16(2)
   -\frac13\left(2\cdot\frac12\right)=4>0.
 \label{eq:b0-product-h}
\end{align}

For two Dirac \(\bm5\)'s, one real \(\bm{10}\), and two complex
\(\bm4\)'s of \(Sp(4)\),
\begin{equation}
 b_0^{Sp(4)}
 =11-\frac43(2)-\frac16(3)
 -\frac13\left(2\cdot\frac12\right)
 =\frac{15}{2}>0.
 \label{eq:b0-sp4}
\end{equation}
For the illustrative \(SU(4)\) list of two Dirac fundamentals, two real
adjoints needed for the breaking, and two complex antifundamentals,
\begin{equation}
 b_0^{SU(4)}
 =\frac{44}{3}-\frac43(1)-\frac16(8)-\frac13(1)
 =\frac{35}{3}>0.
 \label{eq:b0-su4}
\end{equation}
All three pass the one-loop asymptotic-freedom screen.  This says nothing by
itself about which Higgs vacuum is nonperturbatively selected.

\section{Classification and exact status}

\begin{longtable}{@{}L{0.17\textwidth}L{0.20\textwidth}L{0.18\textwidth}L{0.35\textwidth}@{}}
\toprule
Deep group & relevant global form & charge screen & conclusion\\
\midrule
\endfirsthead
\toprule
Deep group & relevant global form & charge screen & conclusion\\
\midrule
\endhead
\(SU(3)\) & unbroken \(H_1\) & fail & charge-one singlet forbidden for every irrep\\
\(SU(2)_c\times SU(2)_H\) & direct product & pass & minimal semisimple backup; opposite-charge partners\\
\(SO(4)\) & diagonal quotient & fail & \((\bm1,\bm2)\) scalar does not descend\\
\(Sp(4)=Spin(5)\) & simply connected & pass & best simple candidate in scanned set\\
\(SO(5)\) & \(Sp(4)/\ZZ_2\) & fail & required spinor \(\bm4\) scalar absent\\
\(SU(4)\) & simply connected & pass & unique fundamental block equality; more elaborate breaking\\
\(SU(4)/\ZZ_2\) & central quotient & fail & \(\bm4,\overline{\bm4}\) do not descend\\
\bottomrule
\end{longtable}

The word ``best'' refers only to the scanned algebraic criteria: a single
simple simply connected group, injective unbroken subgroup, minimal
representations containing both charges, and positive one-loop \(b_0\).
It is not a claim that \(Sp(4)\) dynamically realizes the desired phase.

\begin{longtable}{@{}L{0.44\textwidth}L{0.17\textwidth}L{0.30\textwidth}@{}}
\toprule
Claim & Status & Boundary\\
\midrule
\endfirsthead
\toprule
Claim & Status & Boundary\\
\midrule
\endhead
Product representatives \(G_3,G_2\) & \status{proved} & transverse inverse images generate both reduced summands\\
Circle and torus mod-two indices & \status{proved} & exact Fourier spectra\\
Reference ambient product phase \(+1\) & \status{proved for stated regulator} & paired four-dimensional spectrum only\\
Four defect-regulator character values & \status{proved} & Pontryagin--Thom localization and KO indices\\
Unique charged-QC2D eta phase & \status{open} & microscopic operator and five-dimensional APS data unspecified\\
Diagonal-quotient obstruction & \status{proved} & \(2j+q=0\pmod2\)\\
Simply connected \(Sp(4)\) charge embedding & \status{proved} & explicit \(\bm4,\bm5,\bm{10}\) branching\\
Exactly-two neutral triplet operator & \status{proved at group level} & needs two complex \(\bm4\) copies; no unit-cell derivation\\
Displayed partner split & \status{tree-level} & \(m=-yV\) is unprotected\\
Displayed anomaly and running screen & \status{passed} & vectorlike/Witten parity and one-loop \(b_0\) only\\
Full \(Sp(4)\) gauge-bordism phase & \status{open} & complete Dai--Freed character not evaluated\\
All-scale \(k=+2\) orientation & \status{open} & heavy threshold eta counterterm unmatched\\
Nonperturbative phase and \(B=1\) soliton & \status{open} & requires new strong-dynamics calculation\\
Degree-one Route-E portal & \status{blocked} & waits for all preceding gates\\
\bottomrule
\end{longtable}

\section{Deterministic regression}

The companion verifier
\begin{center}
\texttt{route\_f/code/verify\_ap\_e3\_aps\_global\_form\_search.py}
\end{center}
performs the following checks without fitted parameters:
\begin{enumerate}
\item counts R/NS circle zero modes and mass-path spectral crossings;
\item counts all four flat-torus spin-structure kernels and verifies the Arf
  parity;
\item computes the minimum ambient product gaps and a truncated eta-pairing
  residual;
\item constructs all four defect-regulator character values;
\item enforces the diagonal-centre descent condition;
\item derives the \(Sp(4)\) \(\bm5\) from \(\wedge^2\bm4-\bm1\) and the
  \(\bm{10}\) from \(\operatorname{Sym}^2\bm4\);
\item checks the \(Spin(5)\) centre image, \(SU(4)\) block kernel, and the
  fundamental \(SU(N)\) uniqueness equation;
\item evaluates the displayed mass splits, Witten parities, light level, and
  one-loop beta coefficients;
\item asserts every fail-closed APS, threshold, dynamics, and portal flag.
\end{enumerate}
The numerical calculations encode and regress the displayed analytic
arguments; they do not independently prove a bordism theorem or strong
dynamics.

\section{Ordered decisive calculations}

Several strategies remain available, in descending priority.
\begin{enumerate}
\item \textbf{Microscopic eta route.}  Write the full charged-QC2D
  \(\mathcal D_{\rm UV}\) of Eq.~\eqref{eq:actual-uv-operator}, including all
  heavy fields and regulator masses.  Construct mapping tori representing
  \(G_3,G_2\), then compute the two exponentiated reduced eta invariants.
  This is the only route that selects \((\epsilon_3,\epsilon_2)\) for that
  UV theory.
\item \textbf{Simply connected \(Sp(4)\) route.}  Derive a symmetry or
  controlled threshold prescription protecting \(M_+=0\), evaluate the
  full gauge-bordism character, and match the heavy determinant across
  \(Sp(4)\to SU(2)_c\times U(1)\).
\item \textbf{Minimal-product backup.}  If the simple-group threshold is
  obstructed, repeat the determinant calculation in direct
  \(SU(2)_c\times SU(2)_H\), retaining the global form and both scalar
  copies explicitly.
\item \textbf{Nonperturbative gate.}  For any surviving candidate, establish
  the charged two-colour vacuum, monopole effects, and the complete \(B=1\)
  Hessian before composing with AP-E4.
\end{enumerate}

Until these steps close,
\begin{equation}
 \boxed{\begin{gathered}
 \texttt{route\_e\_uv\_completion\_closed=false},\\
 \texttt{physics\_promotion\_allowed=false},\\
 \texttt{degree\_one\_portal\_allowed=false}.
 \end{gathered}}
 \label{eq:final-gates}
\end{equation}

\bibliographystyle{unsrt}
\bibliography{route_f/tex/ap_e3_aps_global_form_search}

\end{document}
